\documentclass[11pt]{article}
\usepackage[utf8]{inputenc}
\usepackage[margin=1in]{geometry}
\usepackage{hyperref}
\usepackage{enumitem}

\title{Prompt Design References for RQ1a Judge Evaluation}
\author{}
\date{November 2025}

\begin{document}

\maketitle

\section{Overview}

This document provides citations for the three prompt variants used in the RQ1a corruption detection experiments. Each prompt type is grounded in established research on LLM-as-a-judge evaluation strategies and causal reasoning.

\section{Prompt Design Variants}

\subsection{Baseline Prompt}

\textbf{Design Philosophy}: Simple, direct evaluation approach asking ``Does this relationship make sense?''

\textbf{Purpose}: Serves as a control condition and simplicity benchmark for comparison with more sophisticated prompting strategies.

\textbf{Theoretical Grounding}:
\begin{itemize}[noitemsep]
    \item Monte Carlo Data (2024) establishes best practices for LLM-as-a-judge systems, recommending clear evaluation criteria and structured output formats~\cite{montecarlo2024}.
    \item EvidentlyAI (2024) provides comprehensive guidance on implementing LLM-based evaluation systems, emphasizing the importance of baseline comparisons~\cite{evidentlyai2024}.
\end{itemize}

\subsection{Mechanistic Prompt}

\textbf{Design Philosophy}: Domain-specific causal reasoning evaluation focusing on mechanistic evidence.

\textbf{Evaluation Focus}: Assesses \textit{how} the source variable affects the target, including:
\begin{itemize}[noitemsep]
    \item Mechanism description (causal pathway)
    \item Causal direction support
    \item Temporal precedence
    \item Logical coherence
\end{itemize}

\textbf{Theoretical Grounding}:
\begin{itemize}[noitemsep]
    \item MDPI Electronics (2024) demonstrates that large language models require specialized prompting strategies for causal reasoning tasks, as they often conflate correlation with causation without explicit guidance~\cite{mdpi2024}.
    \item TechRxiv (2024) introduces the Logic-of-Thought framework, which emphasizes verification-based reasoning for causal inference tasks~\cite{techrxiv2024}.
    \item ArXiv (2024) presents PC-SubQ (Prompting with Causal Subquestions), showing that fixed, domain-specific subquestions outperform general prompting strategies for enabling LLMs to infer causation from correlation~\cite{arxiv2024pcsubq}.
\end{itemize}

\subsection{SCE (Structured Criteria Evaluation) Prompt}

\textbf{Design Philosophy}: Multi-dimensional structured evaluation with explicit, independent criteria.

\textbf{Evaluation Framework}: Four-dimensional assessment:
\begin{enumerate}[noitemsep]
    \item \textbf{Mechanism Specificity}: Does the explanation describe \textit{how} the causal relationship operates?
    \item \textbf{Causal Direction}: Is the direction of influence (source $\rightarrow$ target) clearly justified?
    \item \textbf{Temporal Order}: Is temporal precedence (cause before effect) addressed?
    \item \textbf{Logical Coherence}: Is the reasoning internally consistent and plausible?
\end{enumerate}

\textbf{Theoretical Grounding}:
\begin{itemize}[noitemsep]
    \item \textbf{Wolfe (2024)} introduces the G-Eval framework, demonstrating that structured criteria-based evaluation significantly outperforms unstructured chain-of-thought approaches for LLM evaluation tasks. This serves as the primary theoretical foundation for SCE~\cite{wolfe2024geval}.
    \item Arize AI (2024) provides evidence-based strategies for LLM-as-a-judge systems, emphasizing the importance of explicit evaluation dimensions~\cite{arizeai2024}.
    \item Monte Carlo Data (2024) establishes that combining explicit criteria with low-precision scoring (e.g., 0.0, 0.5, 1.0) and structured output formats improves reliability and interpretability~\cite{montecarlo2024}.
    \item Medium/Data Science Collective (2024) analyzes when to use reasoning steps versus explanations, finding that structured verification questions are more effective than verbose chain-of-thought for evaluation tasks~\cite{medium2024}.
\end{itemize}

\section{Key Research Insights}

\begin{enumerate}
    \item \textbf{Structured criteria outperform unstructured CoT}: G-Eval (Wolfe, 2024) demonstrates that explicit evaluation dimensions yield more reliable and interpretable judgments than asking LLMs to ``think step by step''~\cite{wolfe2024geval}.
    
    \item \textbf{Domain-specific prompts improve causal reasoning}: PC-SubQ (ArXiv, 2024) shows that fixed, causal-specific subquestions significantly improve LLM performance on causal inference tasks compared to general prompting~\cite{arxiv2024pcsubq}.
    
    \item \textbf{Verification beats generation}: Logic-of-Thought framework (TechRxiv, 2024) demonstrates that verification-based reasoning (checking specific criteria) is more reliable than generation-based reasoning for complex inference tasks~\cite{techrxiv2024}.
    
    \item \textbf{Best practices for LLM judges}: Research consensus (Monte Carlo, 2024; EvidentlyAI, 2024) recommends explicit criteria, structured outputs, and low-precision scoring for robust LLM-as-a-judge systems~\cite{montecarlo2024,evidentlyai2024}.
\end{enumerate}

\section{Methodological Justification}

The three-prompt design tests sensitivity of judge performance across a spectrum of prompting strategies:

\begin{itemize}
    \item \textbf{Baseline}: Minimal scaffolding (control condition)
    \item \textbf{Mechanistic}: Domain-specific guidance (causal reasoning focus)
    \item \textbf{SCE}: Structured multi-criteria evaluation (maximum rigor)
\end{itemize}

This design allows systematic investigation of whether judge agreement and accuracy improve with increasingly sophisticated prompting strategies, grounded in established LLM evaluation research.

\begin{thebibliography}{9}

\bibitem{arizeai2024}
Arize AI.
\textit{Evidence-Based Prompting Strategies for LLM-as-a-Judge}.
2024.
\url{https://arize.com/blog-course/llm-as-judge-prompting-strategies/}

\bibitem{arxiv2024pcsubq}
ArXiv.
\textit{PC-SubQ: Prompting Strategies for Enabling Large Language Models to Infer Causation from Correlation}.
2024.
\url{https://arxiv.org/abs/2410.01420}

\bibitem{evidentlyai2024}
EvidentlyAI.
\textit{LLM-as-a-Judge: A Complete Guide}.
2024.
\url{https://www.evidentlyai.com/llm-as-a-judge}

\bibitem{mdpi2024}
MDPI Electronics.
\textit{Evaluating Causal Reasoning Capabilities of Large Language Models}.
2024.
\url{https://www.mdpi.com/2079-9292/13/18/3618}

\bibitem{medium2024}
Medium/Data Science Collective.
\textit{LLM-as-a-Judge: When to Use Reasoning, Chain-of-Thought, and Explanations}.
2024.
\url{https://medium.com/data-science-collective/}

\bibitem{montecarlo2024}
Monte Carlo Data.
\textit{LLM-As-Judge: 7 Best Practices for Building Robust Evaluation Systems}.
2024.
\url{https://www.montecarlodata.com/blog-llm-as-judge-best-practices/}

\bibitem{techrxiv2024}
TechRxiv.
\textit{Logic-of-Thought Framework for Causal Inference with Large Language Models}.
2024.
\url{https://www.techrxiv.org/}

\bibitem{wolfe2024geval}
Wolfe, C.
\textit{Using Large Language Models for Evaluation: The G-Eval Framework}.
2024.
\url{https://cameronrwolfe.substack.com/p/using-llms-for-evaluation}

\end{thebibliography}

\end{document}
